\paragraph{Smaller finite fields using roots of two.}
When the field also contains a $d$th root of two, the same decoder can
use the factors $2^i-1$ instead of $2^{di}-1$. This reduces the integer
product bound substantially. The following statement requires both
specified roots in the field; it does not replace the prime-progression
hypothesis of Corollary~\ref{ou:asymptotic}.

\begin{proposition}[A root-of-two refinement]\label{ou:kummer}
Let $d$ be prime and let $\omega,\alpha\in\F_p$ satisfy
$\operatorname{ord}(\omega)=d$ and $\alpha^d=2$.
Take $m\ge2$, $1\le D<m$, $0\le c<d$, and put
\[
 n=d(m+1)+c,\qquad K=dD-1,\qquad
 A=n2^{\lceil(m+c+1)/d\rceil}.
\]
Assume
\[
 p>A^{d(d-1)},\qquad p>2^{D(2m-D+3)/2}.
\]
Use the core $\{\alpha^i\omega^j:2\le i\le m+1,\ 0\le j<d\}$,
extra coordinates $\alpha^{m+2},\ldots,\alpha^{m+c+1}$, and padding
$\{\omega^j:0\le j<d\}$.
These points are distinct. The corresponding locator line has exactly
$\binom mD$ nearby parameters at radius $1-K/n-(d+1)/n$;
each has a unique codeword, recoverable in polynomial time in $\log p$.
Parameter zero has exact distance $1-K/n-1/n$.
The construction and decoder are deterministic given the two roots.
\end{proposition}
\begin{proof}
Let $u=2^{1/d}>0$ and let $\zeta$ be a complex primitive $d$th root.
The field $L=\mathbb Q(u,\zeta)$ has degree $d(d-1)$: the degrees
$d$ and $d-1$ of its two subfields are relatively prime, so their
product both divides and bounds $[L:\mathbb Q]$.
The ring $\mathbb Z[u,\zeta]$ is free on the basis
$u^s\zeta^j$ for $0\le s<d$, $0\le j<d-1$.
The given finite-field roots define a ring homomorphism to $\F_p$.
Consequently a nonzero element mapping to zero has nonzero integer
norm divisible by $p$, as follows by reducing its multiplication
matrix modulo $p$.

For a lifted subset sum or a difference of two displayed formal
coordinates, every conjugate has absolute value at most $A$.
Thus its norm has absolute value at most $A^{d(d-1)}<p$.
A modular zero must therefore already be zero in $L$. Distinct formal
coordinates have different positive magnitudes unless their indices
agree, in which case their root-of-unity factors are distinct. This
proves distinctness after reduction.

For a subset sum of nonpadding coordinates, group exponents of $u$
according to their residue modulo $d$, using $u^{qd+s}=2^qu^s$.
The powers $1,u,\ldots,u^{d-1}$ are linearly independent over
$\mathbb Q(\zeta)$, so each resulting polynomial in $\zeta$ vanishes.
Its degree is at most $d-1$, forcing all its coefficients to agree.
Within each residue class, each coefficient is a sum of distinct
powers of two. Binary uniqueness again forces whole core orbits and
excludes all extra coordinates.

The all-codeword classification from Theorem~\ref{ou:finite} now
applies, with locators
\[
 H_I(X)=\prod_{i\in I}(X^d-2^i),\qquad
 I\subseteq\{2,\ldots,m+1\},\quad |I|=D.
\]
Their signed labels lift to positive integers
$T_I=\prod_{i\in I}(2^i-1)<2^{D(2m-D+3)/2}<p$.
Writing $s=\sum_{i\in I}i$, the normalized product satisfies
\[
 \tfrac12<\prod_{i\in I}(1-2^{-i})<1,
\]
so $T_I$ has exactly $s$ binary digits. At index $i$, membership is
equivalent to the remaining normalized product being at most
$1-2^{-i}$; the product of all later factors is strictly greater than
that threshold. The greedy decoder therefore works in base two.
Starting the indices at two is essential to this normalization argument.
The remaining degree, count, and distance conclusions are unchanged.
\end{proof}

For a Mersenne prime $p=2^b-1$ with $d\ne b$, one can take
$\alpha=2^v$, where $dv\equiv1\pmod b$. The check
$\alpha^d=2$ is immediate. An order-$d$ root is still required, so
$d\mid p-1$ remains a separate condition.

Exact certificates give the following specific domains. Rates are $K/n$
and need not be one half. Each comparison uses that row's exact binary
entropy, evaluated through integer powers rather than floating-point logs.
\begin{center}
\begin{tabular}{@{}rrrrll@{}}
\toprule
Prime & $d$ & $n$ & $K$ & Separation & Ratio $>$\\
\midrule
$2^{521}-1$&2&80&29&$2\eta/3$&$2^3$\\
$2^{1279}-1$&3&198&68&$3\eta/4$&$2^4$\\
$2^{9689}-1$&5&905&324&$5\eta/6$&$2^{14}$\\
$2^{23209}-1$&13&3588&1338&$13\eta/14$&$2^2$\\
$2^{44497}-1$&19&7068&2830&$19\eta/20$&$1$\\
\bottomrule
\end{tabular}
\end{center}
The ratio is $\binom mD/(n2^{H_2(K/n)/\eta})$, with
$\eta=(d+1)/n$. The respective $(m,D,c)$ are
$(39,15,0)$, $(65,23,0)$, $(180,65,0)$, $(275,103,0)$,
and $(371,149,0)$.
All five primes have independent Lucas--Lehmer replays; saved roots,
domain distinctness, the two field bounds, strict Elias, and representative
witness evaluations are checked independently. These are deterministic
specific-domain certificates once the certified roots are supplied.
The final row has a uniquely recoverable nearby witness at every accepted
parameter, even though the far point is $19/20$ of the capacity gap
outside the radius and the numerical prescription is too small.
